\documentclass[UTF8, a4paper, 10.5pt, twocolumn, fontset=none]{ctexart}

\usepackage[top=2.0cm, bottom=2.0cm, left=1.8cm, right=1.8cm, columnsep=0.6cm]{geometry}
\setlength{\columnseprule}{0pt}

\usepackage{fontspec}
\usepackage{xeCJK}
\setCJKmainfont[BoldFont={PingFang SC Semibold}, ItalicFont={Kaiti SC}]{Songti SC}
\setCJKsansfont{PingFang SC}
\setCJKmonofont{STSong}
\xeCJKsetup{CJKecglue={\hskip 0.15em plus 0.05em minus 0.05em}}
\usepackage{booktabs}
\usepackage{tikz}
\usetikzlibrary{positioning, arrows.meta, shapes.geometric, fit, backgrounds}
\usepackage[table]{xcolor}
\usepackage{graphicx}
\usepackage{hyperref}
\usepackage{xurl}
\usepackage{caption}
\usepackage{float}
\usepackage{titlesec}
\usepackage[sort&compress]{gbt7714}

\pagestyle{plain}

\newcommand{\rev}[1]{#1}
\setlength{\parindent}{2em}
\linespread{1.18}
\captionsetup{font=small, labelsep=quad}
\titleformat{\section}{\large\bfseries}{\thesection}{0.6em}{}
\titleformat{\subsection}{\normalsize\bfseries}{\thesubsection}{0.5em}{}
\titleformat{\subsubsection}{\normalsize\bfseries}{\thesubsubsection}{0.5em}{}

\sloppy

\hypersetup{
  hidelinks,
  pdfborder={0 0 0},
}

\title{Playwright 测试隐性误修复研究\\
\large \mdseries Research on silent false healing in Playwright test repair}

\author{
  高东彪\\[-1pt]
  \small（软视软件（苏州）有限公司，江苏 苏州 \,215000）\\[2pt]
  \small GAO Dongbiao\\[-1pt]
  \small (Zoom Video Communications (Suzhou), Inc., Suzhou 215000, China)
}
\date{}

\begin{document}

\twocolumn[
  \begin{@twocolumnfalse}
  \maketitle
  \vspace{-12pt}
  \begin{quote}
  \small
  \noindent\textbf{作者简介：}高东彪（1986—），男，工程师，学士；研究方向：软件测试。\\
  \noindent\textbf{*通信作者：}高东彪（1986—），男，工程师，学士；研究方向：软件测试。

  \vspace{4pt}
  \noindent\textbf{摘\ 要：}针对 Playwright 自动化测试中定位器易失效的问题，文章研究大语言模型驱动测试自愈的隐性误修复风险。该研究以 ReproBreak 数据集 koenig 项目的真实失效用例为对象，构建 HealReact 静态基底，并在选择器解析层面分析修复行为。结果表明，当目标元素缺失时，朴素修复器易生成指向错误元素的可解析选择器，而静态基底可覆盖多数真实失效目标。文章认为，持续集成场景应引入行为重放验证机制，以降低误修复风险。

  \vspace{2pt}
  \noindent\textbf{关键词：}测试自愈；大语言模型；Playwright；React；隐性误修复

  \vspace{2pt}
  \noindent\textbf{中图分类号：}TP311.5\quad\textbf{文献标识码：}A

  \vspace{4pt}
  \noindent\textbf{Abstract：} This paper investigates the risk of silent false healing in large language model (LLM)-driven self-healing for Playwright tests with brittle locators. Using real break cases from the koenig project in the ReproBreak dataset, the study builds HealReact, an abstract syntax tree-based static substrate, and analyzes repair behavior at the selector-resolution level. The results show that, when the target element is missing, a naive repairer easily generates resolver-valid selectors for wrong elements, while the static substrate covers most real break targets. The study indicates that continuous integration scenarios require a behavioral replay verification mechanism to reduce false-healing risk.

  \vspace{2pt}
  \noindent\textbf{Key words：} self-healing test; large language model; Playwright; React; silent false healing

  \vspace{6pt}
  \end{quote}
  \end{@twocolumnfalse}
]

\setcounter{section}{-1}
\section{引言}
\label{sec:intro}

脆弱用户界面（User Interface，UI）测试套件的维护是现代 Web 工程中的持续性成本来源。一次层叠样式表（Cascading Style Sheets，CSS）类名重命名、一次 React 组件重新嵌套，或一次行为不变的处理函数重构，都足以使数十个端到端（End-to-End，E2E）Playwright 测试失效。这一问题在图形用户界面（Graphical User Interface，GUI）测试录制重放与自动生成研究综述中被多次列为关键挑战~\cite{li2022androidReplay, wang2025mobileGui}。国内行业期刊也持续关注软件测试自动化方法及测试技术应用形势，相关研究为工程化测试维护问题提供了应用背景~\cite{yang2017wirelessAutoTest, liu2022wirelessTesting}。无论是工业界工具（如 Healenium~\cite{healenium}）还是持续增长的学术工作~\cite{joseph2026zerocost, mouraReproBreak2026}，相关方案通常都将\emph{自愈}作为应对路径，即在定位器失效时由修复器自动将其重写为指向新文档对象模型（Document Object Model，DOM）中等价元素的形式。近年来，大语言模型（Large Language Model，LLM）被广泛引入软件测试与修复任务~\cite{xiang2025llmRepair, li2025llmFuzzing}，自动测试修复（含定位器修复）已成为其中较直接的受益方向。

这一研究方向虽然前景乐观，但现有文献容易给人以问题已接近解决的印象。Joseph~\cite{joseph2026zerocost} 通过一个确定性可访问性阶梯在单一演示样例上报告了 100\% 通过率；当代基于 LLM 的定位器修复工作通常借助一致性安全约束报告修复成功率或解释一致性。然而，这些指标共享一处结构性盲区：相关研究衡量的是修复后的测试是否通过，而不是补丁是否真正指向开发者期望的元素。对于总能找到“某个”可解析选择器的修复器而言，这种盲区具有隐蔽风险，因为每一次“通过但错指”的修复都可能掩盖一次本应暴露为回归的真实 UI 变更。

文章在 ReproBreak~\cite{mouraReproBreak2026} 数据集的真实 Playwright 失效用例上直接测量了这一失败模式。在 koenig 的 75 个可达失效用例上，该研究得到以下结果：

\begin{itemize}
\item \textbf{第一项发现：隐性误修复是默认行为，且软提示\rev{放弃约束}无效。}当真实目标元素被强制从候选集中移除时，朴素 LLM 修复器（\texttt{qwen2.5-coder:7b}~\cite{qwenCoderTech2024}）在 \textbf{78.7\%} 的用例中给出语法正确但元素错误的选择器（59/75）；仅 17.3\% 的情形下能正确放弃。在系统提示中追加明确的“如无候选契合则放弃”\rev{约束}后，结果数字\emph{完全一致}（两版本均为 59/75，$\Delta=0$），表明模型并未将该规则纳入决策。
\item \textbf{第二项发现：两个低开销确定性手段带来零回归的局部改进。}HealReact 的 L1 静态基底在 koenig 上可达 80.6\%（75/93）的真实失效目标；在可达用例上，L3 修复器\rev{首位候选选择器}的精确匹配率为 77.3\%（58/75）。进一步地，testId 加权检索与强锚点后置过滤经二因子析因分离后各带来边际改进、合计 +3 个新正确用例且零回归（方向一致，但在 $n=75$ 下 McNemar 未达显著）。
\item \textbf{第三项发现：跨应用泛化受\rev{锚点约定}差异的制约。}在另一 React+Playwright 代码库（\texttt{payloadcms/payload}，319 个失效配对）上，由于其采用 BEM 风格的模板字面量类名、当前静态抽取器无法对其求值，原始可达率从 koenig 的 80.6\% 跌至 3.8\%。本文将其诊断为一处明确且可修复的盲区，而非对方法本身的否证。
\end{itemize}

\noindent 文章的前两个发现为无人值守自愈流水线引入行为重放验证层提供了直接的经验性动机。文章的第三个发现给出一条可证伪的跨应用泛化论断：任何仅在单一应用上完成基准测试的定位器修复工具，都可能高估其泛化能力，因为锚点约定（\texttt{data-testid}、\texttt{data-kg-*} 与 BEM 模板字面量）在不同 React 代码库之间存在显著差异。

此外，文章构建了 HealReact 的 L1 基底。该基底由抽象语法树（Abstract Syntax Tree，AST）抽取器和带确定性后置校准的 LLM 派生意图标注器组成，在静态层面可达 koenig 93 个失效目标中的 80.6\%。在这一基底之上，小型本地 L3 修复器在可达用例上的正确选择器比例为 77.3\%（在完整数据集上的静态修复代理指标为 62.4\%\footnote{文章全文使用“静态修复代理指标”而非“端到端修复率”：该指标是选择器解析层面的替代评估指标。解析器在元素层面的精确一致性与 Playwright 通过率相关，但既不是其下界，也不是其上界。静态精确匹配的选择器可能因可见性或时序约束在运行时失败，静态不匹配的选择器也可能在重渲染 DOM 中运行通过。两者之间的换算仍需基于 Docker 的真实重放。}）。testId 加权检索与强锚点后置过滤这两个低开销确定性杠杆经二因子析因分离后分别带来边际改进，合计新增 3 个正确用例且未引入回归。这一结果表明，确定性后处理可作为比提示工程更可靠的安全网。

\noindent\textbf{文章贡献：}（i）在真实 Playwright 数据集上对隐性误修复的诊断性测量，以及对软提示放弃指令无效的展示（\S\ref{sec:findings}）；
（ii）HealReact 的 L1 基底作为一个可证伪、可本地运行的静态层（\S\ref{sec:l1}）；
（iii）两个确定性安全杠杆经二因子析因分离，各自带来边际改进、合计 3 增 0 减（方向一致但在 $n=75$ 下 McNemar 未显著，见 \S\ref{sec:findings}）。

该研究测得的 78.7\% 是选择器层面的对抗性上界。若将该数值与 koenig 中真正不可达的 18/93 用例作朴素相乘，则可得到无人值守工作负载下约 15\% 的粗略上界估计，而非实测部署率，详见 \S\ref{sec:findings}。即便仅作为上界估计，文章在 \S\ref{sec:discussion} 中仍指出，该风险已经高到不宜在持续集成（Continuous Integration，CI）场景下无人值守地自动提交 LLM 提出的选择器补丁。

\section{背景与相关工作}
\label{sec:background}

\subsection{定位器脆性与基准数据集}
ReproBreak~\cite{mouraReproBreak2026} 是迄今规模最大的可复现定位器失效数据集，其中包含来自 359 个带 Cypress/Playwright 测试 Web 仓库的 449 个可复现失效。文章将 \texttt{tryghost/koenig}~\cite{koenig} 的 Playwright 切片作为主要研究基底，该切片包含 93 个可复现失效，分布于 19 个版本提交。

\subsection{自愈工具}
测试自愈可视为测试失败后的自动修复问题，即系统需要在既有测试语义约束下生成候选补丁，并避免接受“通过但错误”的补丁。自动程序修复综述为这一修复框架提供了相邻背景~\cite{li2019programRepair}。Healenium~\cite{healenium} 在运行时通过 DOM 相似度修复 Selenium 选择器，是常见的工业界基线。Joseph~\cite{joseph2026zerocost} 基于 DOM 可访问性树确定性地重新抽取选择器，在单个电子商务演示样例上报告了 100\% 通过率，是确定性静态修复轴上的代表性已发表基线，也是与本文 L1 基底最直接的比较对象。按照文章的检索范围，尚未见有同行评议的 LLM 定位器修复工具在真实 Playwright 基准数据集上报告隐性误修复率；与之最接近的当代工作主要关注修复成功率或解释一致性，两者都不直接约束本文探针所针对的失败模式。

\subsection{缺失的部分}
上述工作均未在真实基准数据集上隔离\emph{隐性误修复}，即修复器给出语法通过但指向错误元素的选择器，并由此遮蔽真实 UI 变更的比率。Joseph 在一个目标答案按构造一定可达的演示样例上报告了通过率；ReproBreak 提供了数据基底，但未在其上评估任何修复器。文章在 \S\ref{sec:findings} 中通过一个对抗性构造的探针填补了这一空缺。

\subsection{React 相关细节}
\texttt{koenig-lexical}~\cite{lexicalRichText, koenig} 是一个基于 React 与 Lexical 的现代富文本编辑器，其 Playwright 套件在很大程度上依赖自定义 \texttt{data-kg-*} 属性（约占其失效的一半），而不是 \texttt{data-testid}。\texttt{payloadcms/payload}~\cite{payloadCMS} 是一个 Next.js 后台 UI，其 CSS 类名遵循块、元素、修饰符（Block Element Modifier，BEM）方法~\cite{bemMethodology}，并使用 \texttt{\$\{baseClass\}\_\_suffix} 模板字面量。这些字面量对 React 可见，但对纯字面量 AST 抽取器不可见。文章对这两种锚点约定的比较构成了跨应用发现的依据（见 \S\ref{sec:f3}）。

\section{HealReact L1 静态基底}
\label{sec:l1}

HealReact 设想为一条按“静态 $\to$ 行为验证”递进的四层流水线（编号即由静态程度递减、行为验证程度递增的次序给出）：\textbf{L1}（静态 AST 锚点基底，本文）从源码静态抽取候选元素；\textbf{L2}（运行时 DOM 捕获层）补足仅在运行时才渲染、静态不可见的元素；\textbf{L3}（LLM 选择器修复器，本文）在前两层产出的候选上推断修复后的选择器；\textbf{L4}（行为重放验证器）以真实重放证伪“通过但错指”的补丁。这套 L1--L4 编号是本文对 HealReact 系统的自定义层级命名，并非引自外部框架。本文\emph{实现并评估} L1 与 L3；L2 与 L4 在本文中仅作设计占位（分别对应运行时 DOM 捕获与行为重放验证，均留作后续工作）。

HealReact 的 L1 层是 L3 修复器与\rev{未来行为验证器（L4）所依托}的静态基底。它由三部分组成（图~\ref{fig:arch}）。

\begin{figure}[t]
\centering
\resizebox{\columnwidth}{!}{%
\begin{tikzpicture}[
  font=\scriptsize\sffamily,
  every node/.style={align=center},
  box/.style={draw, rounded corners=2pt, inner sep=3.5pt, minimum height=7mm, minimum width=18mm, fill=white},
  layer/.style={draw, rounded corners=3pt, inner sep=5pt, fill=blue!12},
  arrow/.style={-{Latex[length=1.6mm]}, thick},
  dashbox/.style={draw, dashed, rounded corners=2pt, inner sep=4pt, fill=gray!8},
]
\node[box] (ast) {AST 抽取器\\\texttt{ts-morph}};
\node[box, right=3mm of ast] (intent) {意图标注器\\\texttt{qwen2.5:3b}};
\node[box, right=3mm of intent] (cal) {校准器\\规则 A--J};
\begin{pgfonlayer}{background}
\node[layer, fit=(ast)(intent)(cal), label={[font=\bfseries, yshift=1mm]above:L1 静态基底（本文）}] (l1) {};
\end{pgfonlayer}

\node[box, below=8mm of intent, minimum width=40mm, fill=green!20] (sheet)
  {\texttt{LocatorSheet.json}\\koenig 上\rev{每提交约 400 条记录}\\80.6\% 可达率};

\node[dashbox, below left=8mm and 4mm of sheet, minimum height=12mm, minimum width=24mm, fill=gray!18] (rt)
  {\textbf{L2} 运行时\\DOM 捕获\\（后续工作）};
\node[box, below=8mm of sheet, minimum height=12mm, minimum width=29mm, fill=orange!20] (heal)
  {\textbf{L3} 修复器\\\texttt{qwen2.5-coder:7b}\\+ 检索 + 强锚点};
\node[dashbox, below right=8mm and 4mm of sheet, minimum height=12mm, minimum width=28mm, fill=gray!18] (oracle)
  {\textbf{L4} 行为重放\\验证器\\（后续工作）};

\draw[arrow] (ast) -- (intent);
\draw[arrow] (intent) -- (cal);
\draw[arrow] (l1.south) -- (sheet.north);
\draw[arrow] (sheet.south) -- (heal.north);
\draw[arrow, dashed] (rt.east) -- node[above, font=\tiny]{运行时补充} (heal.west);
\draw[arrow] (heal.east) -- node[above, font=\scriptsize]{补丁} (oracle.west);
\draw[arrow] (oracle.north) |- node[pos=0.25, right, font=\tiny]{接受/拒绝} (sheet.east);
\end{tikzpicture}%
}
\caption{HealReact 流水线。L1（本文）是静态 AST 锚点基底，意图标注层在本文中仅在 32 元素的小规模\rev{构造样例集}上评估；L3 是 \S\ref{sec:findings} 中评估的 LLM 修复器。L2（运行时 DOM 捕获）与 L4（\rev{行为重放验证器}）以虚线框表示，二者在本文中均为设计占位、留作后续工作：L2 负责静态不可见的运行时元素，L4 的必要性由 \S\ref{sec:findings} 中的对抗性误修复探针在经验上得到充分动机化。}
\label{fig:arch}
\end{figure}

\subsection{AST 抽取器}
基于 \texttt{ts-morph}~\cite{tsmorph} 构建的抽取器遍历目标目录下的每个 \texttt{.tsx}/\texttt{.jsx} 文件，并为每个交互式 \rev{JavaScript XML（JSX）}元素发出一条 \rev{JavaScript 对象表示法（JavaScript Object Notation，JSON）}\texttt{LocatorSheet} 记录。每条记录包含元素 tag、父链、role、\texttt{data-testid}、\texttt{aria-label}、\texttt{id}、\texttt{className}、自由文本内容，关键还包括\emph{任意}自定义 \texttt{data-*} 属性（捕获在 \texttt{dataAttrs} 映射中）。抽取器还把任何承载稳定锚点（\texttt{data-testid}、\texttt{data-cy}、\texttt{data-kg-*}、\texttt{aria-label}，或驼峰 \texttt{dataTestId}/\texttt{testId}/\texttt{testID} prop）的 JSX 元素都视为可交互 —— 而不仅是规范\rev{超文本标记语言（HyperText Markup Language，HTML）}tag 集。仅这一项改动就将 \texttt{koenig-lexical} 上的锚点覆盖率从 22\% 提升到 48\%、解析器命中率从 0\% 提升到 71\%，说明现代 React 的脆性中有相当部分藏在自定义包装组件里。

\subsection{意图标注器与校准器}
HealReact 还包含一个 LLM 意图标注流程（经 Ollama~\cite{ollama} 运行的 \texttt{qwen2.5:3b}），它依据组件源码与 JSX 上下文为每个交互元素生成简短的功能性意图标签（例如 \texttt{submit-order-button}、\texttt{toggle-email-notification}）；随后由一个十条规则（A--J）构成的确定性校准层负责识别并降低可疑标签的置信度，可疑模式包括兄弟集合的塌缩标签、把 HTML tag 直接\rev{当作}名词的退化、以及 role 与 intent 不匹配等。被标记的记录将被重新提交给 \texttt{qwen2.5-coder:7b} 重打标签，形成“先低开销、后高开销”的分级验证机制。\textbf{范围说明}：本文对意图层的评估限定在 32 元素的\rev{构造样例集}上；\S\ref{sec:findings} 中针对真实失效用例的所有数字仅使用抽取器输出的锚点与词法检索，\emph{不}使用 L1 意图标签。真实代码库上对意图稳定性的大规模评估超出本文范围。

\subsection{定位器解析器作为可证伪的透镜}
文章构建了一个 Playwright 选择器解析/匹配器，它接受任意选择器表达式（\texttt{page.locator}、\texttt{getByTestId}、\texttt{getByRole}、\texttt{waitForSelector}、原始 CSS）并返回它会匹配的 LocatorSheet 记录集。这是 \S\ref{sec:findings} 中各项论断的\rev{判定工具}：可达性是“\rev{失效}选择器是否能解析到 sheet 中的某条记录？”；L3 正确性是“解析器是否把 LLM 提议的选择器与\rev{真值选择器}映射到同一元素？”。构建这个\rev{判定工具}要求研究处理引号感知的嵌套参数、动态 \rev{JavaScript（JS）}表达式属性值、以及裸属性选择器；早期版本由于\rev{解析器缺陷}把可达率虚抬到 85--98\%，文章现在在诚实性审计中明确标注了这一点。

\subsection{核心指标}
在 \texttt{koenig-lexical} 的 19 个失效\rev{版本提交}上，抽取器在每个\rev{版本提交}中平均产出约 400 条 LocatorSheet 记录（\rev{跨提交}范围 314--489），解析器在静态层面可达真实 Playwright 失效目标的 \textbf{75 / 93 = 80.6\%}。剩余 18 个用例需要运行时 DOM（例如 Lexical 框架在运行时才渲染出的 \texttt{<p>} 节点），属于经典的运行时 DOM 捕获范畴，由 HealReact 的 L2 层负责。

\section{三个实证发现}
\label{sec:findings}

下文结果均来自同一套实验脚本、数据记录与统一判定规则，相关过程可重复复算。该实验覆盖两个 L3 变体和两个误修复探针。

\subsection{复现配置}
为保证可比性，文章固定了以下设置。L3 修复器与意图标注器分别经 Ollama 运行 \texttt{qwen2.5-coder:7b} 与 \texttt{qwen2.5:3b}，两者均取 \texttt{temperature}=0。候选检索采用加性词元重叠打分，权重为 testId$\times 3$、dataAttr key/value$\times 2$、aria/id/class/text/path$\times 1$，并从排序结果中截取前 $k$ 个候选。对于每个失效用例，实验均记录其失效选择器、真值选择器、模型给出的候选、解析结果以及是否判为误修复或精确匹配。\textbf{代码与数据可用性：}实验脚本、结果记录与论文快照已公开归档于 \url{https://github.com/vgawen/healreact}，表~\ref{tab:falseheal}--\ref{tab:cross-app} 中的计数结果均可由该归档复算。

\subsection{共享 L3 协议}
在 \S\ref{sec:f1} 与 \S\ref{sec:f2} 中，L3 修复器均经由 Ollama 以 \texttt{temperature}=0 调用 \texttt{qwen2.5-coder:7b}。用户提示包含以下三部分：(i) 原样给出的失效 Playwright 选择器；(ii) 失效行附近的 $\pm 3$ 行测试源码；(iii) 来自 LocatorSheet 的前 $k$ 个候选，并按加性词元重叠打分检索（testId $\times 3$、dataAttr key/value $\times 2$、aria/id/class/text/path $\times 1$）。模型必须给出三行输出：\texttt{SELECTOR}、\texttt{CANDIDATE\_IDX} 与 \texttt{RATIONALE}。文章未采用 JSON 输出格式，因为模型曾多次在选择器中输出未转义引号。对于多匹配情况，解析器按 \texttt{(componentFile, line)} 顺序取第一个匹配记录作为解析目标；多匹配选择器会被标记，但指标计算中仍使用第一个命中。

\subsection{目标缺失下的隐性误修复}
\label{sec:f1}

\subsubsection{设置}
对 75 个\emph{可达}的 koenig 失效用例（即失效版本提交处 LocatorSheet 中存在真值元素的用例），文章在相同提示与检索条件下重跑 L3 修复器，但将真值元素从 LocatorSheet 与候选集中强制移除，使模型在字面意义上无法选中正确答案。此时，正确响应应为放弃（\texttt{CANDIDATE\_IDX: -1}）；任何给出可解析选择器的响应都会构成一次隐性误修复，并沉默地将测试改写到错误元素上。

文章比较了两个提示变体：
\begin{enumerate}
\item \textsc{Vanilla}：与 L3 基线相同的系统提示。
\item \textsc{Abstain}\rev{：}同一提示外加一条明确的“如无候选契合则放弃”\rev{约束}，告诉模型 \emph{“一次沉默地把测试改写到错误元素上的盲修复永远比放弃更糟”}。
\end{enumerate}

\subsubsection{结果}
表~\ref{tab:falseheal} 给出结果。

\begin{table}[t]
\centering
\caption{误修复探针（\S\ref{sec:f1}）。两个提示变体产出\emph{完全一致}的数字\rev{：}软提示\rev{放弃约束}并不降低隐性误修复。}
\label{tab:falseheal}
\small
\resizebox{\columnwidth}{!}{%
\begin{tabular}{lrrrr}
\toprule
\rev{提示变体} & 放弃\,(好) & 误修复\,(坏) & 未解析 & \textbf{误修复\,\%} \\
\midrule
\textsc{Vanilla} & 13 & 59 & 3 & \textbf{78.7\%} \\
\textsc{+ \rev{放弃约束}} & 13 & 59 & 3 & \textbf{78.7\%} \\
\bottomrule
\end{tabular}}
\\[2pt]
{\scriptsize $n=75$ 个可达 koenig 失效用例。真值目标已从候选集中移除；任何被给出的可解析选择器即为一次隐性误修复。78.7\%（59/75）的 Wilson 95\% 置信区间约为 [68.1\%, 86.4\%]。}
\end{table}

为排除“聚合计数偶然相同”的可能，文章在表~\ref{tab:falseheal-paired} 中给出两提示变体的逐用例配对交叉表。

\begin{table}[t]
\centering
\caption{Vanilla 与 Abstain 的逐用例配对交叉表（\S\ref{sec:f1}，$n=75$）。所有非对角元素均为 0，表明两变体在每一个用例上给出\emph{完全相同}的判定，而非仅聚合计数相同。}
\label{tab:falseheal-paired}
\small
\begin{tabular}{lrrr}
\toprule
Vanilla＼Abstain & 放弃 & 误修复 & 未解析 \\
\midrule
放弃\,(好)   & 13 & 0 & 0 \\
误修复\,(坏) & 0  & 59 & 0 \\
未解析       & 0  & 0 & 3 \\
\bottomrule
\end{tabular}
\\[2pt]
{\scriptsize 在 \texttt{temperature}=0 的确定性解码下，75 个用例的逐例标签两两一致（\rev{非对角项} $=0$，$\Delta=0$）。}
\end{table}

\subsubsection{解读}
该研究得到的 78.7\% 是隐性误修复的上界，因为每个用例都属于对抗性构造。在混合工作负载下，只有 L1 sheet 缺失目标的用例面临此风险，本实验中该比例为 koenig 失效用例中的 18/93 = 19\%。若作朴素相乘，则可得到 $\approx 0.19 \times 0.79 = 15\%$ 的粗略风险估计，而不是实测部署率，因为对抗性探针的行为不必一对一迁移到真正仅在运行时可达的用例上。即便仅作为风险估计，文章仍认为这一风险已经高到不宜允许 CI 在无人值守条件下自动提交 LLM 提议的选择器补丁。

\textsc{Vanilla} 与 \textsc{Abstain} 之间的等同性更具诊断意义，而且这一等同性在逐用例层面成立。表~\ref{tab:falseheal-paired} 显示，75 个用例的配对判定无一发生改变。这一结果说明，7B 规模的代码模型在该任务上不会因为软提示式的放弃约束而改变其可观察行为。任何稳健的放弃机制都必须置于模型之外，例如确定性的检索分数阈值、锚点类别归属规则，或能够证伪通过补丁的行为验证器。该结果从实证层面支持将 HealReact 的 L4 层（行为重放验证器）作为下一阶段构建重点，而不是继续依赖提示级别的安全机制。

\subsection{L3 基线与两个低开销确定性手段}
\label{sec:f2}

\subsubsection{设置}
在同样的 75 个可达 koenig 失效用例上，文章评估了两个 L3 变体：
\begin{itemize}
\item \textsc{v0（基线，两杠杆均关）}：候选按与\rev{失效}选择器的词元重叠程度检索，所有锚点（含 testId）统一权重；无后置过滤；系统提示要求模型优先选择 testId 类锚点。
\item \textsc{v1（两杠杆均开）}\rev{：}(i) 检索打分器把 testId 匹配加权（$\times 3$）排在 dataAttr key/value（$\times 2$）和 aria/id/class/text/path（$\times 1$）之上；(ii) 一个确定性强锚点后置过滤器\rev{：}当所选候选携带可静态求值的 testId 但模型给出的不是 testId 选择器时，将其确定性地重写为 \texttt{getByTestId(testId)}；若 testId 来自不可求值的 JSX 表达式，则不做重写。系统提示不变。
\end{itemize}

为分离两个杠杆各自的贡献，文章进一步以二因子析因（testId 加权 $\in\{$关,开$\}$ $\times$ 后置过滤 $\in\{$关,开$\}$）跑满 4 个实验臂（表~\ref{tab:l3-factorial}）。所有实验臂均使用 \texttt{qwen2.5-coder:7b}、\texttt{temperature=0} 和相同系统提示；真值是 ReproBreak 中 \texttt{new\_locator} 所匹配的目标元素（经解析器解析）。

\subsubsection{结果}
表~\ref{tab:l3-levers} 汇总 L3 v0 与 v1 的端到端对比，表~\ref{tab:l3-factorial} 给出 testId 加权检索与强锚点后置过滤的二因子析因结果。

\begin{table}[H]
\centering
\caption{L3 v0 $\to$ v1 端到端对比（75 个可达 koenig 失效）。v1 = v0 + testId 加权检索 + 强锚点后置过滤；v0 为两杠杆均关的统一权重基线。}
\label{tab:l3-levers}
\small
\resizebox{\columnwidth}{!}{%
\begin{tabular}{lrr}
\toprule
指标 & v0（两杠杆关） & v1（两杠杆开） \\
\midrule
\rev{首位候选精确元素匹配} & 55 / 75 (73.3\%) & \textbf{58 / 75 (77.3\%)} \\
有效 Playwright 选择器 & 68 / 75 (90.7\%) & 68 / 75 (90.7\%) \\
静态修复代理指标（全量 93） & 55 / 93 (59.1\%) & \textbf{58 / 93 (62.4\%)} \\
\midrule
强锚点后置过滤触发 & 0 & 1 / 75（正确） \\
新增正确用例（vs v0） & --- & 3 \\
回归用例（vs v0） & --- & 0 \\
\bottomrule
\end{tabular}}
\end{table}

\begin{table}[H]
\centering
\caption{两杠杆二因子析因（\rev{首位候选精确元素匹配}，$n=75$）。行=testId 加权，列=强锚点后置过滤。每个杠杆单独带来 +2、两者合计 +3、零回归；方向一致但在该样本量下未达统计显著。}
\label{tab:l3-factorial}
\small
\begin{tabular}{lrr}
\toprule
 & 过滤：关 & 过滤：开 \\
\midrule
加权：关 & 55 / 75 & 57 / 75 \\
加权：开 & 57 / 75 & \textbf{58 / 75} \\
\bottomrule
\end{tabular}
\\[2pt]
{\scriptsize McNemar 精确检验（双侧）：单开加权或单开过滤相对“两关”均为 $+2/-0$（$p=0.50$）；两者全开相对“两关”为 $+3/-0$（$p=0.25$）。所有差异方向一致、零回归，但在 $n=75$ 下均未达 $p<0.05$。}
\end{table}

\subsubsection{解读}
这两项手段均属于低开销改动，包括一次正则优先级调整和一段约 10 行的后置过滤，而且完全确定、无需模型再训练。借助表~\ref{tab:l3-factorial} 的二因子析因，文章可以分离两者的贡献。以“两杠杆均关”为基线（55/75），单开 testId 加权检索与单开强锚点过滤均带来 +2 个精确匹配，两者全开时合计带来 +3 个精确匹配，且未出现回归。强锚点后置过滤器在过滤单开臂触发 2 次且均正确，在 v1 全开臂触发 1 次且同样正确，从而消除了部分“LLM 选中正确候选但在生成选择器时挂到不稳定锚点”的失败模式。需要如实说明的是，这些增益虽然方向一致、且未造成回归，但在 $n=75$ 的样本量下经 McNemar 精确检验均未达统计显著（$p\ge0.25$）。因此，文章将其定位为“低开销、零回归的确定性工程杠杆”，而不是已经被统计显著性确立的性能提升。这一结果也与 \S\ref{sec:f1} 的结论相互呼应，即在确定性介入有效之处，提示级介入的作用相对有限。

\subsection{跨应用泛化与 BEM 盲区}
\label{sec:f3}

\subsubsection{设置}
任何仅在单一应用上完成的实验都可能受到应用特定约定的影响。为此，文章补充选择 \texttt{payloadcms/payload}~\cite{payloadCMS} 作为第二个 React+Playwright 代码库进行对照分析。文章在其 \texttt{packages/ui/src} 上运行未修改的 L1 抽取器，产出 457 条记录，并对 319 个失效配对使用 \texttt{new\_locator} 字符串进行解析匹配，以观察跨应用可达性的变化。

\subsubsection{结果}
\begin{table}[t]
\centering
\caption{跨应用 L1 可达率（\S\ref{sec:f3}）。原始下降是真实的，但成因可诊断。}
\label{tab:cross-app}
\small
\resizebox{\columnwidth}{!}{%
\begin{tabular}{lrr}
\toprule
指标 & \texttt{koenig} & \texttt{payload} \\
\midrule
抽取器记录数（均值/单次） & \rev{每提交约 400} & 457 \\
Playwright 失效配对 & 93（已复现） & 319（CSV） \\
\texttt{new\_locator} 可达 & 75 / 93 = \textbf{80.6\%} & 12 / 319 = \textbf{3.8\%} \\
\midrule
主导锚点 & \texttt{data-kg-*}（约 50\%） & CSS class$^\dagger$ \\
模板字面量 class? & 罕见 & 普遍（BEM） \\
\bottomrule
\end{tabular}}
\\[2pt]
{\scriptsize $^\dagger$payload 的主导锚点占比来自对 307 个未命中用例的采样，而非全量失效分布。}
\end{table}

原始 3.8\% 的可达率较低，但文章对 307 个未命中用例的检查结果呈现出清晰规律。payload 使用 BEM~\cite{bemMethodology}：每个组件声明一个 \texttt{baseClass} 常量（例如 \texttt{const baseClass = 'dashboard'}），并以 \texttt{\{`\$\{baseClass\}\_\_add-button`\}} 这样的模板字面量写出 className。测试引用的是渲染后的类 \texttt{.dashboard\_\_add-button}，而纯字面量 AST 抽取器无法直接看到这一结果。

\subsubsection{解读}
这一结果并非对 HealReact L1 设计的否证，而是识别出一处范围明确的盲区，其修复对应一次确定性的 AST 处理：构造 \texttt{baseClass} 解析器，回溯常量声明并展开模板字面量，并配以 CSS Modules~\cite{cssModules} 等价的展开器。基于对 307 个未命中用例的抽样观察，文章推测修复后原始可达率会明显抬升，但这一判断仍需进一步实测验证。当前结果表明，锚点约定的多样性（\texttt{data-testid}、自定义 \texttt{data-*}、BEM 模板字面量等）很可能是定位器修复工具跨应用泛化的重要约束，任何仅在单一应用上完成基准测试的工具都可能高估其泛化能力。

\section{讨论与局限}
\label{sec:discussion}

\subsection{78.7\% 对无人值守 CI 自愈意味着什么}
LLM 自愈常见的部署模式具有较高自主性：一次端到端测试失败会触发 LLM 生成补丁，随后由 CI 自动提交补丁，或将其作为拉取请求（Pull Request，PR）评论发出。根据 \S\ref{sec:f1} 的对抗性上界作粗略换算，在不带行为验证器的此类部署中，约 15\% 的真实 Playwright 失效存在被朴素修复器静默改写到错误元素的风险，但这一比例只是粗略上界估计，并非实测部署率。该风险的危险并不在于“修复器让构建失败”，因为构建失败本就是测试应承担的职责，而在于“修复器使构建通过，但通过的却是错误目标”。

\subsection{为什么软提示安全网会失效}
\textsc{Vanilla} 与 \textsc{Abstain} 两种提示之间的等同性（表~\ref{tab:falseheal}）具有明确启示。一种可能的解释是，系统提示中“必要时放弃”的软指令需要与“生成合理输出”这一更强先验竞争，因此在多数情况下处于劣势；这一解释属于机理层面的推测，文章未进一步开展训练动力学验证。对工具设计而言，这一结果说明反误修复的安全保障必须置于模型之外，例如确定性的检索分数阈值、锚点类别归属规则，以及能够证伪通过补丁的行为验证器。

\subsection{为什么两个低开销确定性杠杆胜过提示工程}
\S\ref{sec:f2} 中的 testId 加权检索与强锚点后置过滤在机制上较为直接。经二因子析因分离后，这两项手段分别带来 +2 个精确匹配，合计带来 +3 个精确匹配，且未出现回归（表~\ref{tab:l3-factorial}）；增益方向虽然一致，但在 $n=75$ 下经 McNemar 精确检验未达统计显著。因此，文章既不声称其具有统计显著性，也不声称该增益可以泛化到当前设置之外，而是将其视为 HealReact 的一条局部设计经验：与依赖系统提示让模型内化约束相比，将锚点优先级约束放入候选池和输出后置过滤更可靠，且不会引入回归。

\subsection{局限}
\textit{模型单一性带来的有效性威胁。}文章全部数字均基于单一开源小型代码模型（\texttt{qwen2.5-coder:7b}）。前沿大模型（如 GPT-4 类）是否能够在隐性误修复维度上显著缩小这一差距，仍需后续实测；根据 \textsc{Vanilla = Abstain} 等同性所揭示的机理，文章倾向于认为，若缺乏外部确定性验证器，单纯依靠模型规模扩大未必足以弥合这一差距。
\textit{基准数据集范围。}文章主结果来自单一 React+Playwright 代码库 koenig；\S\ref{sec:f3} 已扩展到 payload，但其可达率仍受到 BEM 类名盲区的限制。
\textit{静态修复代理指标与运行时通过率。}77.3\% 与 62.4\% 这两个 L3 数字反映的是选择器解析器层面的精确一致性。其与真实 Playwright 运行时通过率之间的换算关系，仍需在 ReproBreak 提供的逐版本提交复现环境中通过 Docker 重放修复后的测试来建立。作为最小可行性检查，文章对 3 个 v1 静态精确匹配用例（编号 615、702、703）进行了 Docker 运行时重放：三者的修复基线均通过；替换为 HealReact 提议选择器后，2 个用例运行通过，1 个用例运行失败（编号 702，等待 \texttt{getByTestId('color-picker-toggle')} 超时）。这一小规模验证实验无法估计全量通过率，但说明静态精确匹配既存在可迁移到运行时的案例，也可能因可见性、时序或上下文差异而失败。因此，文章继续将 77.3\% 与 62.4\% 标为静态修复代理指标，而不宣称其等同于端到端通过率。

\subsection{有效性威胁}
78.7\% 这一数字依赖解析器语义。具体而言，\texttt{resolve\_locators.py} 将动态 JS 表达式属性值视为可能匹配，并对复合 CSS 使用祖先过近似；这两项决策在错误排查场景下具有合理性，但在边界情况下偏松。文章在解析器假设审计中已撤回早期“可达率上限 98\%”的虚高声明，并在收紧解析器假设后将 80.6\% 作为保守数字。由于解析器是文章所有“正确匹配、误修复、可达率”判定的事实判定器，其判定边界会直接影响这些数字，因此文章进行了 strict/permissive 双假设的敏感性扫描（表~\ref{tab:resolver-sens}）。文章主要实验结果中的 80.6\%（75/93）来自 permissive 假设；若关闭上述两处过近似并改用 strict 假设，可达率将降至 54.8\%（51/93）。因此，L1 静态可达率的合理理解区间应为 [54.8\%, 80.6\%]，其中 80.6\% 对应上界。此外，在 permissive 假设下，75 个可达用例中有 46 个为多匹配，解析器按 \texttt{(componentFile, line)} 顺序取首个命中消歧；未解析用例数则从 permissive 下的 18 个增至 strict 下的 42 个。对于抽样用例，文章尚未使用真实 DOM 快照对解析判定进行交叉校验，这部分工作仍留待后续完成。

\begin{table}[t]
\centering
\caption{解析器 strict/permissive 敏感性（koenig，$n=93$）。permissive 为本文所采用的假设；strict 关闭“动态属性值视为可能匹配”与“同文件祖先近似”两处过近似。}
\label{tab:resolver-sens}
\small
\resizebox{\columnwidth}{!}{%
\begin{tabular}{lrr}
\toprule
指标 & permissive（本文） & strict（下界） \\
\midrule
new\_locator 可达 & 75 / 93 (80.6\%) & 51 / 93 (54.8\%) \\
多匹配（$\ge2$ 命中） & 46 & 0 \\
未解析（0 命中） & 18 & 42 \\
\bottomrule
\end{tabular}}
\end{table}

\subsection{未来工作}
文章的诊断结论延伸出两类可推进的工程化方向：其一，将静态基底向跨锚点约定（BEM 模板字面量、CSS Modules 等）扩展，以提高不同 React 项目中的可达率稳定性；其二，基于 \S\ref{sec:findings} 给出的经验性必要性，构建并集成一个基于网络行为重放的轻量级验证器，作为 LLM 修复结果的下游安全机制。

\section{结论}
\label{sec:conclusion}

在缺乏行为验证器的部署条件下，LLM 驱动的 Playwright 自愈与“测试通过但真实回归被遮蔽”之间仅一步之遥。文章在对抗性极限下于选择器解析层面测得 78.7\% 的隐性误修复上界，而提示级放弃指令并未带来实质缓解；这一数字属于静态诊断指标，其与运行时测试通过率之间的换算仍需基于 Docker 的真实重放。HealReact 的 L1 基底使 koenig 数据集中 80.6\% 的真实失效在静态层面可达；小型本地修复器在可达用例上的首位候选选择器精确匹配率达到 77.3\%；两个低开销确定性手段进一步带来帕累托改进。上述结果为后续评估行为重放验证器提供了直接的经验性动机，也提出了一条受锚点约定差异限制的、可证伪的跨应用泛化假设。


\bibliography{references}

\end{document}
